% A270_Z3_Sektoren_En_ch.tex — English parallel version of A270
% Z3 sector structure and hadron mass correction in T4/Z3. Johann Pascher, 23 July 2026

\section{Purpose and Status}

\textbf{[STIPULATION]} This document investigates the sector structure
of the T$^4$/Z$_3$ orbifold with regard to a structural distinction
between elementary particles (leptons) and composite particles
(hadrons). Central result: the three Z$_3$ sectors assign to each
particle class a topologically fixed $K_\text{frak}$ exponent ---
leptons $0$, mesons $+1$, baryons $+2$ --- in addition to a common bulk
exponent~$36$.

\textbf{Status:} hypothesis at consistency level [K]. The meson anchor
is passed via the pion binding (A155). The baryon step is a prediction;
its direct test requires a geometric baryon base mass (open).

\section{Starting Point: Bulk Exponent 36}

\subsection{The Key Relation}

\textbf{[K]} Two independently determined FFGFT quantities meet in one
number. First, the D4 packing density $\pi^2/16$, whose reciprocal is
\[
  \frac{16}{\pi^2} \approx 1{.}62114.
\]
Second, the scale-recursion depth from Doc.~275,
\[
  k^* = \frac{\ln\varphi}{\xi} \approx 3609{.}1
  \quad\Rightarrow\quad \frac{k^*}{100} = 36{.}09.
\]
Numerically:
\[
  K_\text{frak}^{-36} = \left(\frac{75}{74}\right)^{36} \approx 1{.}62130
  \approx \frac{16}{\pi^2} \quad (\Delta = 0{.}010\,\%).
\]

\subsection{Statistical Safeguard}

\textbf{[K]} Monte-Carlo check (100\,000 trials): the hit probability
$\leq 0{.}010\,\%$ for a random target in the value range is
$1{.}56\,\%$ per candidate. Of 21 tested physical constants ($\pi$,
$e$, $\varphi$, $\sqrt2$, $\sqrt3$, $\ln 2$, \ldots) \emph{only}
$16/\pi^2$ hits at $\leq 0{.}010\,\%$; the next-best candidate lies at
$0{.}066\,\%$. The independent support by $k^*/100$ (Doc.~275 was
determined without reference to masses) considerably weakens the chance
hypothesis. P35 remains noted: a forward derivation of the identity
$K_\text{frak}^{-36} = 16/\pi^2$ does not exist.

\textbf{[STIPULATION]} The exponent~36 is henceforth called the
\emph{bulk exponent}: the common fractal correction of the untwisted
sector in which all particles live.

\section{T$^4$/Z$_3$ Orbifold Structure}

\subsection{Fixed-Point Geometry}

\textbf{[B]} Z$_3$ acts on T$^3 = S^1_x \times S^1_y \times S^1_z$ by
cyclic permutation:
\[
  (x,y,z) \;\xrightarrow{g}\; (y,z,x) \;\xrightarrow{g}\; (z,x,y)
  \;\xrightarrow{g}\; (x,y,z)
\]
Fixed points: $x = y = z = k\cdot L/3$, $k = 0,1,2$ --- exactly
\textbf{3 fixed points}, isomorphic to Z$_3$.

Local geometry at each fixed point:
\begin{center}
\renewcommand{\arraystretch}{1.3}
{\small\begin{tabular}{lll}
\toprule
Direction & Vector & Z$_3$ sector \\
\midrule
$\mathbf{e}_0 = (1,1,1)/\sqrt{3}$   & invariant   & sector 0 (untwisted) \\
$\mathbf{e}_1 = (1,-1,0)/\sqrt{2}$  & phase $\omega$   & sector $g$ (twisted) \\
$\mathbf{e}_2 = (1,1,-2)/\sqrt{6}$  & phase $\omega^2$ & sector $g^2$ (twisted) \\
\bottomrule
\end{tabular}}
\renewcommand{\arraystretch}{1.0}
\end{center}
At each fixed point: 1 untwisted $+$ 2 twisted degrees of freedom ---
exactly the Z$_3$ sector structure \{$e, g, g^2$\}.

\subsection{Quark Localization and Confinement}

\textbf{[H]} In FFGFT, $\tilde{T}\cdot m = 1$ holds; the mass $m$ sits
on $S^1_m$, one of the three spatial circles. Under Z$_3$ permutation,
$S^1_m$ has two images on the other circles.

\textbf{Elementary particle (lepton):} sits on $S^1_m$ alone,
transforms as a 1-dimensional Z$_3$ representation, remains in the
untwisted sector~0.

\textbf{Composite particle of 3 constituents (baryon):}
constituent~1 on $S^1_x$, constituent~2 on $S^1_y$, constituent~3 on
$S^1_z$. The localization condition ``all 3 at the same spatial
point'' reads:
\[
  x = y = z \quad \Leftrightarrow \quad \textbf{fixed-point condition}
\]
Quark confinement in T$^4$/Z$_3$ is thereby \emph{topologically
enforced}: free quarks would violate the fixed-point condition.
Leptons (1 circle) have no fixed-point condition and are free.

\subsection{Sector Exponents}

\textbf{[H]} $K_\text{frak}$ is a property of the sector partition
function. Each of the 3 Z$_3$ sectors contributes $K_\text{frak}^1$.
The bulk exponent~36 covers the untwisted sector (already contained in
the ladder formulas). Twisted sectors contribute an additional
$K_\text{frak}^1$ each. Hence the sector table:

\begin{center}
\renewcommand{\arraystretch}{1.3}
{\small\begin{tabular}{lccc}
\toprule
Particle class & Const. & Active sectors & Exponent \\
\midrule
Lepton ($e$, $\mu$, $\tau$) & 1 & $\{e\}$ & $36+0$ \\
Meson ($q\bar{q}$) & 2 & $\{e, g\}$ & $36+1 = 37$ \\
Baryon ($qqq$) & 3 & $\{e, g, g^2\}$ & $36+2 = 38$ \\
\bottomrule
\end{tabular}}
\renewcommand{\arraystretch}{1.0}
\end{center}

The exponents are topologically fixed: Z$_3 = \{e,g,g^2\}$ has exactly
two non-trivial elements, independently of any mass.

\section{Empirical Anchor: the Pion}

\textbf{[K]} A155 introduces the additive formula
$m_M = m_{q_1} + m_{q_2} + \Lambda_\text{QCD}\cdot
K_\text{frak}^{n_\text{eff}}$ for mesons. Evaluating the pion binding
yields (A155):
\[
  \boxed{n_\text{eff}(\pi^0) = 37{.}79}
  \qquad n_\text{eff}(\pi^\pm) = 36{.}62.
\]
Both values lie within the window of the sector prediction $36+1 = 37$
for mesons. \textbf{The meson step of the sector table is thereby
empirically anchored at the pion} --- the exponent was determined in
A155 (from the legacy-corpus formula) independently of this sector
argument.

The domain of validity of the additive formula (only binding
$< \Lambda_\text{QCD}$; kaon, $\rho$, $\omega$ structurally outside)
and the GMOR extension are treated in A155.

\section{Consistency Evidence}

\subsection{Koide Relation}

\textbf{[K]} The Koide relation
\[
  Q = \frac{m_e+m_\mu+m_\tau}{(\sqrt{m_e}+\sqrt{m_\mu}+\sqrt{m_\tau})^2}
  = \frac{2}{3}
\]
holds exactly for leptons (A110: the geometric ladder predicts $Q$
within $0{.}16\,\%$ of $2/3$ without using Koide; the experimental
value hits $2/3$ exactly with $\Delta \approx 3\times10^{-5}$).

For quarks, according to A150, \emph{no} corresponding structure
exists:
\begin{quote}\small
``Second, there are six quarks instead of three leptons, but no
ordering corresponding to the $Z_3$ structure that would map them onto
three eigenvalues of one operator. The circulant plane of A110 has no
counterpart here.'' (A150, \S{}Quarks)
\end{quote}

This is consistent with the sector hypothesis --- and sharpens it: the
Z$_3$ circulant (A110), which carries the Koide structure of the
leptons, \emph{is} the algebra of the untwisted sector. Leptons in the
untwisted sector $\to$ circulant acts $\to$ Koide holds. Quarks couple
to twisted sectors $\to$ the circulant ordering has ``no counterpart''
(A150) $\to$ no Koide analogue exists. The sector hypothesis thus
supplies a candidate for the \emph{reason} of the structural deficit
recorded in A150.

\subsection{$K_\text{frak}$ Distribution in the Corpus}

\textbf{[K]} $K_\text{frak}$ appears in the corpus base formulas
exactly where the sector table demands it:

\begin{center}
\renewcommand{\arraystretch}{1.3}
{\small\begin{tabular}{lll}
\toprule
Class & $K_\text{frak}$ in base formula? & Sector table \\
\midrule
Leptons & no ($m_L = r\,\xi^{p}v$, A100) & untwisted: no extra \\
Mesons & yes ($\Lambda\cdot K_\text{frak}^{n}$, A155) & $+1$ sector \\
Baryons & yes; corrected legacy calculation with anchor $\pi^2/2$ (Doc.~190, R61; $\Delta=0{.}05\,\%$) & $+2$ sectors \\
\bottomrule
\end{tabular}}
\renewcommand{\arraystretch}{1.0}
\end{center}

\section{Prediction: Baryon Step}

\textbf{[H]} From the sector table, the exponent for baryons is
$36+2 = 38$:
\[
  K_\text{frak}^{-38} = 1{.}6654.
\]
The corresponding statement: a geometric baryon base mass (before
sector correction) relates to the observed one by
$K_\text{frak}^{-38}$; equivalently, the baryon exponent lies exactly
one unit above the meson value measured at the pion.

\textbf{Direct test open:} it requires an independently derived
geometric baryon base mass, which does not exist in the corpus (A150:
the quark sector is descriptive classification, no predictive
quantity). Until then, the baryon step rests on the topology (two
twisted sectors) and the pion anchoring of the meson step.

\section{Open Points and Falsification}

\begin{enumerate}
\item \textbf{Forward derivation of $K_\text{frak}^{-36} = 16/\pi^2$}
  (P35): the identity is supported numerically and via $k^*/100$, but
  not derived.
\item \textbf{Baryon sector}: the legacy calculation is corrected by
  register entry (Doc.~190, R61). A fully geometric candidate for the
  proton mass ---
  $m_p = (\pi^3/12)\,N_c(1+\alpha_s)\,e^{-3\xi/4}\,\Lambda_\text{QCD}/2$,
  every factor interpreted, $\Delta=+0{.}21\,\%$ absorbable --- is kept
  in A155 \S{}Baryon Candidate ([S]; forward derivation open). The
  direct test of the $+2$ step awaits this derivation.
\item \textbf{$\pi^\pm$/$\pi^0$ difference} ($36{.}62$ vs.\ $37{.}79$):
  does the EM self-energy carry a sector contribution of its own?
  (A155)
\end{enumerate}

\textbf{Falsification hierarchy:}
\begin{enumerate}
\item Meson exponent $36+1$: confirmed at the pion ($37{.}79$, A155);
  extension to the kaon via the GMOR form in A155 as [S] candidate.
\item Koide analogue for quarks: must \emph{not} hold --- consistent
  with A150 (``circulant plane of A110 has no counterpart'')
  \checkmark
\item Baryon exponent $36+2$: falsifiable as soon as a geometric
  baryon base mass exists.
\end{enumerate}

\section{Notation}

Symbols used throughout the A-series are explained in A010.
Specific to this document:

\begin{center}
\renewcommand{\arraystretch}{1.3}
{\small\begin{tabular}{lp{9cm}}
\toprule
Symbol & Meaning \\
\midrule
$Z_3=\{e,g,g^2\}$ & cyclic group; $g$: permutation $(x,y,z)\to(y,z,x)$ \\
$\omega=e^{2\pi i/3}$ & primitive third root of unity (twist phase) \\
sectors $0,g,g^2$ & untwisted sector; the two twisted sectors
(eigenvalues $1,\omega,\omega^2$ of the permutation) \\
$\mathbf{e}_0,\mathbf{e}_1,\mathbf{e}_2$ & eigenbasis at the fixed
point: invariant direction $(1,1,1)/\sqrt3$; twisted directions \\
fixed point & point with $x=y=z$ on $T^3$; three of them ($k\cdot L/3$) \\
bulk exponent 36 & common $K_\text{frak}$ exponent of the untwisted
sector; $K_\text{frak}^{-36}\approx 16/\pi^2$ \\
sector exponent & addition per particle class: lepton $+0$, meson
$+1$, baryon $+2$ \\
$k^*=\ln\varphi/\xi$ & scale-recursion depth to $1/\varphi$
(Doc.~275); $k^*/100=36{.}09$ \\
$\pi^2/16$ & D4 packing density; reciprocal $16/\pi^2$ \\
$n_\text{eff}$ & empirical binding exponent from A155 (pion anchor) \\
$Q$ & Koide quantity
$(m_e{+}m_\mu{+}m_\tau)/(\sqrt{m_e}{+}\sqrt{m_\mu}{+}\sqrt{m_\tau})^2$ \\
\bottomrule
\end{tabular}}
\renewcommand{\arraystretch}{1.0}
\end{center}

\medskip\noindent\textbf{Verification script:} \texttt{python/A\_Serie\_Skripte/a270\_z3\_sektoren.py}
--- all numerical statements of this document are checked there with assertions.

\section{References}

\begin{center}
{\small\begin{tabular}{ll}
\toprule
Document & Reference \\
\midrule
A015 & origin of $\xi$, Casimir confirmation \\
A100 & lepton ladder $(r_i, p_i)$ \\
A110 & circulant and Koide: $Q=2/3$, Z$_3$ algebra of the leptons \\
A150 & quark ladder, lepton/quark distinction \\
A155 & meson formula: domain, pion anchor, GMOR; baryon candidate \\
Doc.~190 & P35, P40; register R61 (baryon correction, legacy corpus) \\
Doc.~275 & $k^*\approx3609$, scale recursion, $1/\varphi$ passage \\
Doc.~230/231/232 & T$^4$/Z$_3$ Hilbert-space bridge \\
\bottomrule
\end{tabular}}
\end{center}
